\section{Geneza i przeznaczenie}

SabberStone to symulator gry Hearthstone o otwartym kodzie źródłowym. Został napisany w języku C\# z wykorzystaniem platformy .NET Core i jest udostępniony na licencji AGPL-3.0 \cite{sabberstone_github}. Projekt powstał w 2016 roku jako odpowiedź na ograniczenia wcześniejszych symulatorów, które słabo radziły sobie z implementacją złożonych interakcji między kartami. Autorzy chcieli dodać mechaniki takie jak aury, wzmocnienia (ang. \textit{buffs}) czy kaskadowe efekty wyzwalaczy (ang. \textit{triggers}). Celem było wierne odwzorowanie mechanik oryginału i jednocześnie udostępnienie czytelnych interfejsów dla badaczy zajmujących się algorytmami decyzyjnymi \cite{sabberstone_hearthsim}.

Tym samym SabberStone zyskał status oficjalnej platformy konkursowej dla \textit{Hearthstone AI Competition}, organizowanego corocznie podczas IEEE Conference on Games \cite{ieee_cog, hearthstone_competition}. Konkurs odbył się trzykrotnie, w latach 2018–2020, każdorazowo zbierając od 30 do 50 zgłoszeń podzielonych między dwie ścieżki \cite{kowalski2023}. Po edycji z 2020 roku konkurs przestał być kontynuowany. Mimo to, wokół projektu wykształciła się aktywna społeczność badawcza, a SabberStone nadal funkcjonuje jako środowisko eksperymentalne w pracach naukowych z obszaru gier karcianych.
